\documentclass[a4paper,12pt]{ctexart}
\usepackage{xcolor}
\usepackage{graphicx}
\usepackage{amsmath}
\usepackage[top=1.8cm, bottom=1.8cm, left=1.8cm, right=1.8cm]{geometry}
\usepackage{array}
\usepackage{hyperref}
\hypersetup{colorlinks=true,linkcolor=blue!70!black}
\linespread{1.35}

\definecolor{defblue}{RGB}{0,80,160}
\definecolor{thinkorange}{RGB}{230,120,20}
\definecolor{formulared}{RGB}{180,30,50}
\definecolor{funboxbg}{RGB}{245,248,252}
\definecolor{examplebg}{RGB}{255,250,240}
\definecolor{practicegreen}{RGB}{40,130,70}

\newenvironment{thinkbox}{
  \vspace{0.3cm}
  \begin{center}
  \begin{minipage}{0.88\textwidth}
  \colorbox{funboxbg}{
  \begin{minipage}{0.94\textwidth}
  \vspace{0.15cm}
  \textbf{\textcolor{thinkorange}{【 想一想 】}}
}{
  \vspace{0.15cm}
  \end{minipage}}
  \end{minipage}
  \end{center}
  \vspace{0.3cm}
}

\newenvironment{definition}[1]{
  \vspace{0.2cm}
  \begin{center}
  \begin{minipage}{0.88\textwidth}
  \fbox{
  \begin{minipage}{0.92\textwidth}
  \vspace{0.1cm}
  \textbf{\textcolor{defblue}{【 #1 】}}
}{
  \vspace{0.1cm}
  \end{minipage}}
  \end{minipage}
  \end{center}
  \vspace{0.2cm}
}

\newenvironment{example}{
  \vspace{0.2cm}
  \begin{center}
  \begin{minipage}{0.88\textwidth}
  \colorbox{examplebg}{
  \begin{minipage}{0.94\textwidth}
  \vspace{0.15cm}
  \textbf{\textcolor{orange!80!black}{【 例题 】}}
}{
  \vspace{0.15cm}
  \end{minipage}}
  \end{minipage}
  \end{center}
  \vspace{0.2cm}
}

\newenvironment{practice}{
  \vspace{0.3cm}
  \begin{center}
  \begin{minipage}{0.88\textwidth}
  \colorbox{funboxbg}{
  \begin{minipage}{0.94\textwidth}
  \vspace{0.15cm}
  \textbf{\textcolor{practicegreen}{【 练一练 】}}
}{
  \vspace{0.15cm}
  \end{minipage}}
  \end{minipage}
  \end{center}
  \vspace{0.2cm}
}

\newenvironment{figplaceholder}[1]{
  \vspace{0.2cm}
  \begin{center}
  \fbox{
  \begin{minipage}{0.82\textwidth}
  \centering
  \textcolor{blue!70!black}{\textbf{【插图位置】}} \\[0.2cm]
  #1
  \end{minipage}}
  \end{center}
  \vspace{0.2cm}
}

\newcommand{\formula}[1]{\begin{center}\textcolor{formulared}{\large\boxed{#1}}\end{center}}

\title{\textcolor{blue!70!black}{\Huge\textbf{物理2 · 声光热}}}
\author{}
\date{}

\begin{document}
\maketitle
\thispagestyle{empty}

\section*{\textcolor{blue!70!black}{致同学}}

在力学入门里，我们学会了描述运动、分析力、计算压强和浮力。现在我们要探索三个全新的世界——\textbf{声音、光和热}。这些都是你每天用耳朵、眼睛和皮肤就能感受到的物理现象。你会发现：物理不止是推导公式，它还能解释为什么音乐会好听、潜水时看到的鱼"不在那个位置"、以及冬天铁栏杆为什么比木头"更冷"。

\newpage

\section{\textcolor{orange!80!black}{第一课\quad 声音的产生与传播}}

\subsection*{声音是怎么来的？}

敲一下音叉，音叉发出了嗡嗡的声音。用手摸音叉的顶端——你在发抖！把正在发声的音叉轻轻碰一下水面——水面出现了细密的水花。

这些现象说明同一个道理：\textbf{声音是由物体的振动产生的。}

\begin{definition}{声音的产生}
声音是由物体的\textbf{振动}产生的。一切正在发声的物体都在振动；振动停止，发声也停止。
\end{definition}

人说话时声带在振动，蜜蜂飞过时翅膀在振动（每秒约200-400次），打鼓时鼓膜在振动——你能想到的任何声音，源头都是某个东西在振动。

\subsection*{声音是怎样传到耳朵里的？}

振动产生了声波。声波需要\textbf{介质}来传播——在空气中，声源的振动让周围的空气分子一会儿挤在一起（密部）、一会儿散开（疏部），这种疏密相间的波动向四面八方传播，就是\textbf{声波}。

\begin{definition}{声音的传播}
声音的传播需要\textbf{介质}。固体、液体、气体都可以传声，但\textbf{真空不能传声}。
\end{definition}

为什么真空不能传声？因为真空中没有分子——再也没有东西可以用来传递振动了。太空中如果发生了一场大爆炸，你在旁边也听不见——虽然火光就在眼前。

\textbf{声音在不同介质中的速度：}固体最快、液体次之、气体最慢。在15℃空气中，声速约为\textbf{340 m/s}。在水里，声速约1500 m/s。在铁轨中，声速可以高达5000 m/s——这就是为什么把耳朵贴在铁轨上可以提前听到远处火车的动静。

\subsection*{回声}

声音碰到障碍物会被反射回来——这就是\textbf{回声}。如果原声和回声到达耳朵的时间间隔大于0.1秒，人耳就能把它们区分开来。根据这个原理，声源至少要离障碍物\textbf{约17米}以上，才能听到清晰的回声（声速340 m/s × 0.1s ÷ 2 = 17m）。

\begin{example}
一个人对着山崖大喊一声，2秒后听到回声。此人离山崖多远？（声速 $v = 340\ \text{m/s}$）

\textbf{解：}声音走了来回 $\rightarrow$ 声音单程时间 $t = 2 \div 2 = 1\ \text{s}$。\\
距离 $s = vt = 340 \times 1 = 340\ \text{m}$。
\end{example}

\begin{thinkbox}
为什么在空房间里说话听起来"空洞"，而放满了家具的房间声音就自然了？因为空房间里声音被平行墙面反复反射产生\textbf{混响}，家具则能把声波打散、吸收。音乐厅的墙壁要做成特殊弧面——就是控制回声的科学。
\end{thinkbox}

\begin{practice}
1. （判断）声音可以在真空中传播。\quad （\quad）\\
2. 敲鼓时，鼓面在\_\_\_\_\_；人说话时，\_\_\_\_\_在振动。\\
3. 在15℃的空气中，声速约为\_\_\_\_\_ m/s。\\
4. （计算）雷声在闪电后4秒才听到。雷暴大约在多远？（声速取340 m/s）\\
5. （选择）在月球上两个宇航员即使很近也无法直接听到对方说话，因为（\quad）\\
\quad A. 月球上没有空气 \quad B. 月球上有引力 \quad C. 月球上太冷 \quad D. 声音在月球上传播得太快\\
6. （填空）声音在\_\_\_\_\_中传播最快，在\_\_\_\_\_中传播最慢。15℃空气中声速\textbf{\_\_\_\_\_} m/s。\\
7. （判断）只要有物体的振动，就一定能听到声音。（\quad）\\
8. 人耳区分原声和回声的最小时间间隔是\_\_\_\_\_s，因此至少离障碍物\_\_\_\_\_m才能听到回声。
\end{practice}

\newpage

\section{\textcolor{orange!80!black}{第二课\quad 声音的特性}}

\subsection*{为什么不同的声音听起来不一样？}

钢琴弹一个"do"，电吉他弹同一个"do"——明明是同一个音高，为什么一听就知道一个是钢琴、一个是吉他？

声音有三个独立"维度"：\textbf{音调、响度、音色}。每个维度对应一个物理量。

\subsection*{音调——声音的"高"和"低"}

男生的声音低沉，女生的声音尖细——这就是\textbf{音调}不同。

\begin{definition}{音调}
音调指声音的高低，由声源振动的\textbf{频率}决定。频率越高，音调越高。
\end{definition}

频率的单位是\textbf{赫兹（Hz）}，1 Hz = 1次振动/秒。人耳能听到的频率范围大约是\textbf{20-20000 Hz}。低于20 Hz的叫\textbf{次声波}（地震、火山爆发会产生），高于20000 Hz的叫\textbf{超声波}（蝙蝠用超声波"看"东西、医院B超也是超声波）。

\subsection*{响度——声音的"大"和"小"}

\begin{definition}{响度}
响度指声音的强弱（大小），由声源振动的\textbf{幅度}决定。振幅越大，响度越大。
\end{definition}

响度还和\textbf{到声源的距离}有关——你离声源越远，声音越小（声波的能量在传播中扩散衰减）。这就是为什么老师上课要用嗓子大声讲——最后一排的同学离得远。

响度的单位是\textbf{分贝（dB）}。树叶沙沙约10 dB，正常交谈约60 dB，摇滚音乐会约110 dB。长期暴露在90 dB以上会损伤听力。

\subsection*{音色——声音的"指纹"}

同样的音调、同样的响度，不同乐器听起来完全不同——这就是\textbf{音色}。

\begin{definition}{音色}
音色由声源振动的\textbf{波形}决定。不同乐器发出的声音波形不同，所以听起来不一样。
\end{definition}

用示波器（一种能把声波变成图形的仪器）看——钢琴的波形是规则但复杂的折线，小提琴的波形是平滑的波浪，噪音的波形则是乱七八糟的无规律锯齿。

\begin{figplaceholder}
此处插入示波器上三种不同乐器的波形对比图（见 figure/requirements/book2-ch1.txt 插图1）
\end{figplaceholder}

\begin{example}
一个音叉每秒振动440次。它的频率是多少Hz？人耳能听到这个声音吗？

\textbf{解：}频率 = 440 Hz。440 Hz在20-20000 Hz之间——人耳可以听到。440 Hz就是钢琴上的标准音"A"（la音）。
\end{example}

\begin{thinkbox}
为什么录音里自己的声音和平时听到的"不一样"？平时你自己说话的声音，一部分是通过头骨和颅腔直接传到内耳的（骨传导），一部分是通过空气传来。录音只保留了空气中的声音——少了骨传导的"低音"，所以你总觉得录音里的自己声音变尖了——其实那才是别人耳朵里你的真实声音。
\end{thinkbox}

\begin{practice}
1. "男低音"和"女高音"中的"低"和"高"指的是声音的\_\_\_\_\_。\\
2. （选择）调节收音机的音量旋钮，改变的是声音的（\quad）\\
\quad A. 音调 \quad B. 响度 \quad C. 音色 \quad D. 频率\\
3. 超声波和次声波人耳都听不到，因为它们的\_\_\_\_\_不在人耳的听觉范围内。\\
4. （填空）狗能听到人类听不到的某些声音，因为狗的听觉频率\_\_\_\_\_比人更宽。\\
5. （计算）蝙蝠靠超声波导航，已知超声波在空气中速度为340m/s，蝙蝠发射后0.02s收到回波。蝙蝠距障碍物多远？\\
6. （选择）调节乐器的弦的松紧可以改变声音的（\quad）\\
\quad A. 响度 \quad B. 音调 \quad C. 音色 \quad D. 振幅\\
7. （填空）频率的单位是\_\_\_\_\_，人耳能听到的频率范围是\_\_\_\_\_到\_\_\_\_\_。\\
8. （判断）音调越高，声音传播得越快。（\quad）
\end{practice}

\newpage

\section{\textcolor{orange!80!black}{第三课\quad 光的直线传播}}

\subsection*{光源——光从哪来？}

太阳、电灯、蜡烛、萤火虫——它们有一个共同的身份：\textbf{光源}。光源是自身能发光的物体。月亮不是光源——它只是反射了太阳光。

\subsection*{光沿直线传播}

光在\textbf{同种均匀介质}中沿\textbf{直线}传播。这是几何光学最基本的规律。手电筒射出的光束是直的；透过树叶缝隙射下来的阳光是直的；你在雾天看见的汽车灯光束也是直的。

因为光沿直线传播，才有了\textbf{影子}——你身体挡住了一条直线的光路，背后就留下了暗区。关于影子的更多故事，你可以在自然科学小故事的"影子为什么跟着我走"一课里重温。

\begin{definition}{光线}
为了便于描述光的传播路径，我们在物理学中用一条带箭头的直线来表示光传播的路径和方向，这样的线叫做\textbf{光线}。光线并不是真实存在的线，它只是一种理想模型。
\end{definition}

\subsection*{光速——宇宙的限速牌}

光在真空中的速度约为\textbf{$3.0 \times 10^8\ \text{m/s}$}（每秒30万公里）——去月球1.3秒就到。光速是宇宙中最快的速度，任何有质量的物体都不能超过它。

光在不同介质中速度不同：真空中最快，空气中略慢（但几乎一样），水中约$2.25 \times 10^8\ \text{m/s}$，玻璃中约$2.0 \times 10^8\ \text{m/s}$。光速的变化是后面讲折射的基础。

\textbf{光年}是距离单位——光在真空中一年走过的距离，约$9.46 \times 10^{12}\ \text{km}$。它不是时间单位。

\subsection*{小孔成像——两千年前的"照相机"}

在一块纸板上扎一个小孔，放在蜡烛和墙壁之间。你会发现——墙壁上出现了蜡烛火焰的\textbf{倒立}影像（火焰尖朝下）。

这就是\textbf{小孔成像}。原理是光沿直线传播：火焰顶端的光穿过小孔投射到墙壁下方，火焰底部的光投射到墙壁上方——上下颠倒了。早在两千多年前，中国的墨子就记载了这个现象——这是人类历史上最早的"光学实验"之一。

\begin{figplaceholder}
此处插入小孔成像原理示意图：蜡烛→小孔→倒立影像（见 figure/requirements/book2-ch1.txt 插图2）
\end{figplaceholder}

\begin{thinkbox}
为什么夏天树荫下的光斑大多是圆的？那不是树叶缝隙的形状——是太阳的"小孔成像"！每一个树叶间隙就是一个小孔，在树荫地面上投射出太阳的圆形倒像。
\end{thinkbox}

\begin{practice}
1. 下列哪个不是光源？\quad A. 太阳 \quad B. 月亮 \quad C. 萤火虫 \quad D. 点燃的蜡烛\\
2. 光在\_\_\_\_\_介质中沿直线传播。\\
3. 光在真空中的速度约为\_\_\_\_\_ m/s。\\
4. （判断）"光年"是时间单位。（\quad）\\
5. （计算）太阳光传到地球约需500秒。太阳到地球的距离大约是多少米？\\
6. （选择）小孔成像中，屏上成的是（\quad）\\
\quad A. 正立放大的虚像 \quad B. 倒立缩小的实像 \quad C. 倒立的实像 \quad D. 正立的虚像\\
7. （填空）光在真空中的速度为\_\_\_\_\_ m/s，光在水中的速度比在空气中\_\_\_\_\_。\\
8. （判断）小孔成像实验中，小孔越小，成的像越清晰。（\quad）
\end{practice}

\newpage

\section{\textcolor{orange!80!black}{第四课\quad 光的反射}}

\subsection*{反射定律——光怎么"弹"回去}

光照到物体表面，一部分光会被弹回去——这就是\textbf{光的反射}。我们能看见不发光的物体（桌子、书本、人脸），就是因为它们把光反射到了我们眼睛里。

\begin{definition}{光的反射定律}
在反射现象中：\\
① 反射光线、入射光线和法线在\textbf{同一平面}内；\\
② 反射光线和入射光线分居\textbf{法线两侧}；\\
③ \textbf{反射角等于入射角}（$\angle r = \angle i$）。
\end{definition}

其中\textbf{法线}是一条假想的线——经过入射点且垂直于反射面。它就像镜面的"公平裁判"：入射角和反射角都以它为准来测量。

\subsection*{镜面反射 vs 漫反射}

镜子能把阳光反射成一个耀眼的光点——因为镜面非常平滑，所有入射光整齐地按照反射定律弹往同一个方向。这叫\textbf{镜面反射}。

但你从各个方向都能看到一页白纸——虽然纸的表面在微观上凹凸不平，每一束光仍然遵守反射定律，但法线方向各不相同，光被弹往了各个方向。这叫\textbf{漫反射}。

\textbf{注意：漫反射不是"不遵守反射定律"——它在每个微观局部严格遵守反射定律，只是整体上看起来散射了。}人眼之所以能从各个角度看到同一个物体，正是漫反射的功劳。

\subsection*{平面镜成像}

当你照镜子的时候，你在镜子里看到了一个"自己"。这个像是怎么形成的？

\begin{definition}{平面镜成像特点}
① 像和物到镜面的距离\textbf{相等}；\\
② 像和物的大小\textbf{相等}；\\
③ 像和物的连线与镜面\textbf{垂直}；\\
④ 平面镜成的是\textbf{虚像}（不能用光屏承接）。
\end{definition}

\textbf{虚像和实像的区别}：虚像是由反射光线的反向延长线相交形成的——光实际上并没有经过那个位置（镜子后面是一片墙），但你的大脑"以为"光是从那里直线射过来的。实像则是光真实汇聚的地方——比如小孔成像中墙壁上的蜡烛倒影，你可以在那儿放一张纸，影像就"接"在上面。

\begin{example}
小明身高1.6m，站在平面镜前0.5m处。镜子里的像离小明多远？像有多高？

\textbf{解：}像到镜面的距离 = 物到镜面的距离 = 0.5 m。\\
$\therefore$ 像到小明的距离 = $0.5 + 0.5 = 1.0\ \text{m}$。\\
像的高度 = 物高 = 1.6 m（平面镜成的像和物大小相等）。
\end{example}

\begin{thinkbox}
你照镜子时，镜子里的你"左右颠倒了"——但为什么上下没有颠倒？答案其实很简单：镜子颠倒的是"前后"方向（把前和后互换了），左右是跟着前后一起被翻过来的。想不通？对镜子举起右手——镜子里的人举起的是他的左手。但对镜子举起上方的拳头——镜子里的人也是向上举。所以本质上，镜子不是"左右反"，是"前后反"。
\end{thinkbox}

\begin{practice}
1. 入射光线与镜面的夹角为30°，则入射角为\_\_\_\_\_°，反射角为\_\_\_\_\_°。\\
2. （选择）我们能从各个方向看到黑板上的字，是因为光发生了（\quad）\\
\quad A. 镜面反射 \quad B. 漫反射 \quad C. 折射 \quad D. 直线传播\\
3. 一个人站在平面镜前2m处，他的像距他\_\_\_\_\_ m。\\
4. 如果在穿衣镜前只看到自己上半身的像，要想看到全身——应该（\quad）\\
\quad A. 走近镜子 \quad B. 远离镜子 \quad C. 换一面更大的镜子 \quad D. 不可能\\
5. （计算）一束入射光线与镜面的夹角为40°，则入射角是\_\_\_\_\_°，反射角是\_\_\_\_\_°。\\
6. （选择）下列属于光的反射现象的是（\quad）\\
\quad A. 海市蜃楼 \quad B. 水中倒影 \quad C. 小孔成像 \quad D. 筷子在水中"折断"\\
7. （填空）镜面反射和漫反射都\_\_\_\_\_（遵循/不遵循）光的反射定律。\\
8. （判断）平面镜成的是实像，可以用光屏承接。（\quad）
\end{practice}

\newpage

\section{\textcolor{orange!80!black}{第五课\quad 光的折射}}

\subsection*{光穿越边界时会"拐弯"}

把一根筷子斜插进一杯水里——从侧面看，筷子在水面处"折"成了两截。这不是筷子真的断了，是光在水和空气的交界面发生了\textbf{折射}。

\begin{definition}{光的折射}
光从一种介质\textbf{斜射}入另一种介质时，传播方向发生改变的现象，叫做\textbf{光的折射}。
\end{definition}

折射的原因是：\textbf{光在不同介质中传播速度不同。}光在空气中快、在水中慢——当它斜着穿过水面时，一侧先进入水中减速，导致整体传播方向弯折。

\subsection*{折射的规律}

\begin{definition}{光的折射规律}
① 折射光线、入射光线和法线在\textbf{同一平面}内；\\
② 光从空气斜射入水或玻璃中时，折射光线\textbf{靠近法线}（折射角 $<$ 入射角）；\\
③ 光从水或玻璃斜射入空气中时，折射光线\textbf{远离法线}（折射角 $>$ 入射角）；\\
④ 入射角增大时，折射角也增大；\\
⑤ 光\textbf{垂直}射入界面时，传播方向\textbf{不变}。
\end{definition}

核心记忆法：\textbf{光在"快"介质中角度大，在"慢"介质中角度小。}空气中光速最快→空气中角度最大。

\subsection*{折射造成的"错觉"}

\textbf{池水变浅：}站在池边看水底——水底看起来比实际浅。因为水底物体发出的光从水进入空气时，远离法线折射——光到了你眼睛里，你的大脑以为光是沿直线的，就在比实物更高的位置"虚构"了一个虚像。

\textbf{海市蜃楼（简单版）：}在炎热的沙漠或柏油路上，地面附近空气被晒得很热——热空气密度小、光速略快。远处的光经过不同密度的空气层不断折射，形成了远方天空的倒影——看起来像一滩水，其实是天空的折射像。

\textbf{水中的鱼"不在那个位置"：}叉鱼的时候，如果你对准眼睛看到的位置叉下去——肯定会叉偏。因为光从鱼身上经过水→空气的折射后"抬高"了鱼的位置。你应该往看到的鱼的下方瞄准。

\begin{figplaceholder}
此处插入光的折射示意图：筷子"折断"+池水变浅的光路图（见 figure/requirements/book2-ch1.txt 插图3）
\end{figplaceholder}

\begin{example}
一束光从空气以45°射入水中（$n_{\text{水}} \approx 1.33$）。折射角是大于45°还是小于45°？为什么？

\textbf{解：}小于45°。光从空气（快介质）进入水（慢介质），折射光线靠近法线——折射角 $<$ 入射角。
\end{example}

\begin{thinkbox}
在游泳池里从水底往上看岸上的人——看到的影像比实际高还是低？想一想光路：从岸上人发出的光经过空气→水，折射角小于入射角——光进入你眼睛的角度变"陡"了。大脑的直线反推会把人拉得更高。所以在水里看岸上的人，会觉得他"站在半空中"。
\end{thinkbox}

\begin{practice}
1. 光从空气斜射入水中，折射角\_\_\_\_\_入射角（填"大于""小于"或"等于"）。\\
2. （选择）下列现象中由光的折射引起的是（\quad）\\
\quad A. 水中倒影 \quad B. 小孔成像 \quad C. 池水"变浅" \quad D. 立竿见影\\
3. 一束光从水中斜射入空气，折射光线将\_\_\_\_\_法线。\\
4. 叉鱼时应该瞄准所看到鱼的位置的\_\_\_\_\_方（填"上"或"下"）。\\
5. （计算）一束光从空气斜射入某种液体中，入射角为60°，折射角为35°。光是从快介质进入慢介质吗？\\
6. （选择）潜水员在水中看到岸上的树，与实际位置相比（\quad）\\
\quad A. 变高了 \quad B. 变矮了 \quad C. 不变 \quad D. 左右颠倒\\
7. （填空）光从空气射入水中时，折射角\_\_\_\_\_入射角；光从水射入空气时，折射角\_\_\_\_\_入射角。\\
8. （判断）光垂直射入界面时，传播方向不变。（\quad）
\end{practice}

\newpage

\section{\textcolor{orange!80!black}{第六课\quad 透镜及其应用}}

\subsection*{两种透镜}

\textbf{凸透镜：}中间厚、边缘薄。对光有\textbf{会聚}作用。放大镜就是凸透镜——阳光透过凸透镜能聚焦成一个灼热的小亮点，可以用来点火。

\textbf{凹透镜：}中间薄、边缘厚。对光有\textbf{发散}作用。近视眼镜的镜片就是凹透镜。

\begin{figplaceholder}
此处插入凸透镜和凹透镜的外形对比+三条特殊光线图（见 figure/requirements/book2-ch1.txt 插图4）
\end{figplaceholder}

\subsection*{透镜的三条特殊光线}

对于\textbf{凸透镜}（初中阶段重点）：
\begin{enumerate}
  \item \textbf{过光心的光线}：传播方向不变。
  \item \textbf{平行于主光轴的光线}：折射后过焦点。
  \item \textbf{过焦点的光线}：折射后平行于主光轴。
\end{enumerate}

掌握这三条光线，你就能画出凸透镜成像的光路图，进而确定像的位置和性质。

\subsection*{凸透镜成像规律}

这是初中物理最重要的规律之一。物体在不同距离，凸透镜成的像截然不同。设物距为 $u$，焦距为 $f$，像距为 $v$：

\begin{center}
\begin{tabular}{|c|c|c|c|c|}
\hline
\textbf{物距 $u$} & \textbf{像的性质} & \textbf{像距 $v$} & \textbf{大小} & \textbf{应用} \\
\hline
$u > 2f$ & 倒立实像 & $f < v < 2f$ & 缩小 & 照相机 \\
$u = 2f$ & 倒立实像 & $v = 2f$ & 等大 & （分界点）\\
$f < u < 2f$ & 倒立实像 & $v > 2f$ & 放大 & 投影仪 \\
$u = f$ & 不成像 & — & — & （分界点）\\
$u < f$ & 正立虚像 & $|v| > u$ & 放大 & 放大镜 \\
\hline
\end{tabular}
\end{center}

\textbf{记忆口诀：}一倍焦距分虚实，二倍焦距分大小。物近像远像变大（实像时），物近像近像变小（虚像时）。

\subsection*{眼睛——人体的光学仪器}

人眼就是一部高级的"自动对焦照相机"。晶状体相当于一个凸透镜，视网膜相当于光屏。晶状体的焦距可以通过\textbf{睫状肌}调节——看远处时焦距变长，看近处时焦距变短。

\textbf{近视眼：}晶状体太凸（或眼球前后径太长），远处物体的像落在视网膜\textbf{前方}——看远模糊。用\textbf{凹透镜}矫正。

\textbf{远视眼：}晶状体太平（或眼球前后径太短），近处物体的像落在视网膜\textbf{后方}——看近模糊。用\textbf{凸透镜}矫正。

\begin{example}
一个放大镜的焦距为10 cm。用它看一枚邮票，看到正立放大的像。邮票离放大镜的距离可能为多少？（多选：5 cm、10 cm、15 cm、20 cm）

\textbf{解：}要成正立放大的虚像，物距必须 $u < f$。$f = 10\ \text{cm}$，所以只有5 cm满足条件。
\end{example}

\begin{thinkbox}
为什么透过一滴水能看到放大的字？因为水滴在纸面上自然呈凸透镜形状——中间厚边缘薄——这就是一个微型放大镜。和课本上的"水滴放大镜"实验道理完全一样。
\end{thinkbox}

\begin{practice}
1. 凸透镜对光有\_\_\_\_\_作用，凹透镜对光有\_\_\_\_\_作用。\\
2. （选择）物体在凸透镜前20cm处，在透镜另一侧的光屏上得到倒立缩小的像，则透镜的焦距（\quad）\\
\quad A. $f > 20$cm \quad B. $f > 10$cm \quad C. $f < 10$cm \quad D. $10\text{cm} < f < 20\text{cm}$\\
3. 照相机镜头相当于一个\_\_\_\_\_透镜，拍照时物体应在镜头的\_\_\_\_\_之外。\\
4. （填空）近视眼用\_\_\_\_\_透镜矫正，因为它能把光\_\_\_\_\_（会聚/发散）。\\
5. （计算）一个凸透镜焦距$f=15\ \text{cm}$，物体放在距透镜40cm处。像的性质是\_\_\_\_\_、\_\_\_\_\_、\_\_\_\_\_像。\\
6. （选择）投影仪利用了凸透镜成（\quad）\\
\quad A. 正立放大的虚像 \quad B. 倒立放大的实像 \quad C. 倒立缩小的实像 \quad D. 正立等大的虚像\\
7. （填空）一倍焦距分\_\_\_\_\_，二倍焦距分\_\_\_\_\_。\\
8. （判断）凸透镜成的实像都是倒立的。（\quad）
\end{practice}

\newpage

\section{\textcolor{orange!80!black}{第七课\quad 温度与物态变化}}

\subsection*{温度——分子运动的"评分"}

在物理1的力学中我们学过热胀冷缩。温度的本质就是\textbf{分子运动剧烈程度的量度}。温度越高，分子的平均动能越大。

温度计利用液体的热胀冷缩原理工作。常用温标有三种：
\begin{itemize}
  \item \textbf{摄氏度（℃）}：在标准大气压下，冰水混合物为0℃，沸水为100℃。
  \item \textbf{华氏度（℉）}：欧美日常使用。0℉约等于-17.8℃。
  \item \textbf{开尔文（K）}：科学研究用。0K（绝对零度）$\approx -273.15$℃，理论上分子完全停止运动。
\end{itemize}

\subsection*{六种物态变化}

物质在固态、液态、气态之间的转换，统称\textbf{物态变化}。一共有六种，分为三对：

\begin{center}
\begin{tabular}{|c|c|c|c|}
\hline
\textbf{过程} & \textbf{方向} & \textbf{吸/放热} & \textbf{生活实例} \\
\hline
熔化 & 固态→液态 & \textbf{吸热} & 冰化成水 \\
凝固 & 液态→固态 & \textbf{放热} & 水结成冰 \\
汽化 & 液态→气态 & \textbf{吸热} & 水烧开、湿衣服晾干 \\
液化 & 气态→液态 & \textbf{放热} & 冬天玻璃上的水珠、雾 \\
升华 & 固态→气态 & \textbf{吸热} & 樟脑丸变小消失、干冰 \\
凝华 & 气态→固态 & \textbf{放热} & 霜的形成、树枝上的"雾凇" \\
\hline
\end{tabular}
\end{center}

\textbf{记忆规律：}由"紧凑"变"松散"（固→液→气）需要吸收能量（吸热）；由"松散"变"紧凑"（气→液→固）需要释放能量（放热）。

\subsection*{晶体与非晶体}

\textbf{晶体}（冰、食盐、金属）有固定的熔点。加热时温度上升到熔点后，在完全熔化之前温度保持\textbf{不变}（吸收的热量全部用来破坏晶体结构，不升温）。

\textbf{非晶体}（玻璃、沥青、蜡）没有固定的熔点。加热时一边熔化一边升温，软化成糊状再变液态。

\subsection*{蒸发与沸腾的区别}

两者都属于汽化，但有本质区别：

\textbf{蒸发}：在任何温度下都能发生，只在液体\textbf{表面}进行，是一种缓慢的汽化。蒸发吸热，所以夏天洒水后感觉凉快、汗液蒸发带走体热。

\textbf{沸腾}：在特定温度（沸点）下发生，在液体的\textbf{表面和内部同时剧烈}进行。沸点随气压变化——气压越低，沸点越低。所以在高原上烧水不到100℃就开了。

\begin{example}
为什么夏天用电风扇吹，人感觉凉快？温度计放在电风扇前面吹，读数会下降吗？

\textbf{解：}电风扇加速了汗液的蒸发——蒸发吸热带走体表的温度，所以人感觉凉快。但温度计上没有汗液可以蒸发，所以温度计读数\textbf{不变}。这说明"感觉凉快"不等于"空气温度降低"。
\end{example}

\begin{thinkbox}
冰箱制冷利用了什么原理？制冷剂（一种低沸点液体）在冰箱内部的蒸发器里\textbf{蒸发}（吸热→冰箱内部变冷），然后到冰箱背面的冷凝器里\textbf{液化}（放热→把热量排到房间）。冰箱不是"制造冷气"，而是"把热量从里面搬运到外面"。
\end{thinkbox}

\begin{practice}
1. 写出下列物态变化的名称：\\
\quad ① 冰变成水：\_\_\_\_\_ \quad ② 雾的形成：\_\_\_\_\_ \quad ③ 霜的形成：\_\_\_\_\_\\
2. 晶体在熔化过程中温度\_\_\_\_\_（填"升高""降低"或"不变"）。\\
3. （选择）下列措施中能加快蒸发的是（\quad）\\
\quad A. 把蔬菜用保鲜膜包好 \quad B. 把湿衣服摊开晾在通风处\\
\quad C. 把酒精灯盖上灯帽 \quad D. 把蔬菜放冰箱冷藏室\\
4. 在海拔高处烧水，水的沸点比在平原\_\_\_\_\_（高/低）。\\
5. （计算）500g的冰从$-10^\circ\text{C}$加热到$0^\circ\text{C}$，吸收多少热量？（$c_{\text{冰}}=2.1\times10^3\ \text{J/(kg·℃)}$）\\
6. （选择）下列物态变化中属于放热的是（\quad）\\
\quad A. 熔化 \quad B. 汽化 \quad C. 升华 \quad D. 液化\\
7. （填空）晶体与非晶体的主要区别是晶体具有固定的\_\_\_\_\_。\\
8. （判断）蒸发在任何温度下都能发生，沸腾只在特定温度发生。（\quad）
\end{practice}

\newpage

\section{\textcolor{orange!80!black}{第八课\quad 内能与热量}}

\subsection*{内能——"藏在里面的能量"}

物体内部所有分子做无规则运动的动能和分子间相互作用的势能之和，叫做\textbf{内能}。一切物体都有内能——即使冰块也有内能（分子还在振动，只是没那么剧烈）。

内能和温度不是一个概念：温度是"平均分子动能的评分"，内能是"所有分子动能+势能的总和"。一杯开水的温度比一池温水高，但一池温水的内能总量远比一杯开水大——因为它的分子数量太多了。

\subsection*{改变内能的两种方式}

\textbf{做功}：搓手取暖（机械能→内能）、打气筒打气后筒身发热（压缩气体做功）、钻木取火。

\textbf{热传递}：太阳晒热水（热辐射）、暖气片给房间供暖（对流）、手摸热水杯（热传导）。热传递的方向永远是\textbf{从高温物体到低温物体}。

\subsection*{热量与比热容}

\begin{definition}{热量}
在热传递过程中，传递的那部分内能叫\textbf{热量}，用 $Q$ 表示，单位：焦耳（J）。热量不是一个物体"含有"的东西——它是传递的量。
\end{definition}

物体吸收的热量和四个因素有关：\textbf{质量}$m$、\textbf{比热容}$c$、\textbf{升高的温度}$\Delta t$。

\formula{Q = cm\Delta t}

其中 $c$ 是\textbf{比热容}——单位质量的某种物质温度升高1℃所吸收的热量。水的比热容非常大（$4.2 \times 10^3\ \text{J/(kg·℃)}$）——这意味着水升温慢、降温也慢。正是因为这个特性，沿海地区冬暖夏凉（海水夏天吸热、冬天放热，像个巨大的恒温器）。

\subsection*{燃料的热值与热机效率}

1kg某种燃料完全燃烧放出的热量，叫这种燃料的\textbf{热值}，用 $q$ 表示。汽油的热值约$4.6 \times 10^7\ \text{J/kg}$。

燃料燃烧放出的总热量中，只有一部分能被热机转化为机械功。热机的效率：

\formula{\eta = \frac{W_{\text{有用}}}{Q_{\text{总}}} \times 100\%}

现代汽油机的效率一般在25\%-35\%左右——剩下的能量大部分变成了废气带走的热量和摩擦损耗。提高热机效率是节能的核心方向。

\begin{example}
质量为2kg的水，从20℃加热到100℃，需要吸收多少热量？（$c_{\text{水}} = 4.2 \times 10^3\ \text{J/(kg·℃)}$）

\textbf{解：}$\Delta t = 100 - 20 = 80\ \text{℃}$。\\
$Q = cm\Delta t = 4200 \times 2 \times 80 = 6.72 \times 10^5\ \text{J}$。\\
相当于一块普通的手机充电宝全部电能——够给一杯水从室温加热到沸腾。
\end{example}

\begin{thinkbox}
为什么内陆地区昼夜温差大（"早穿棉袄午穿纱"），而沿海地区温差小？因为泥土和沙石的比热容远小于水——白天太阳晒，内陆升温快；晚上散热，内陆降温也快。海水则像一个巨大的"热缓冲器"：白天吸热慢（海边凉快），晚上放热慢（海边温暖）。比热容解释了气候——物理和地理在这里交汇了。
\end{thinkbox}

\begin{practice}
1. 改变内能的两种方式是\_\_\_\_\_和\_\_\_\_\_。\\
2. 质量为1kg的水温度升高1℃吸收的热量是\_\_\_\_\_ J。\\
3. （计算）500g的铝块（$c_{\text{铝}} = 0.88 \times 10^3\ \text{J/(kg·℃)}$）从25℃加热到75℃，吸收多少热量？\\
4. （判断）热量是物体内部"含有"的一种能量。（\quad）\\
5. （填空）热传递的方向总是从\_\_\_\_\_物体到\_\_\_\_\_物体。\\
6. （计算）$2\ \text{kg}$的煤油从$30^\circ\text{C}$加热到$70^\circ\text{C}$，吸收$1.68\times10^5\ \text{J}$热量。煤油的比热容多大？\\
7. （选择）下列现象中利用做功改变内能的是（\quad）\\
\quad A. 晒太阳感到暖和 \quad B. 用热水袋暖手 \quad C. 钻木取火 \quad D. 用冰袋给发烧病人降温\\
8. （填空）热值是燃料的一种特性，汽油的热值约\_\_\_\_\_ J/kg。\\
9. （判断）物体温度越高，所含的热量越多。（\quad）
\end{practice}

\vspace{0.5cm}

\begin{center}
\rule{0.7\textwidth}{1pt}

\vspace{0.4cm}
{\large\textcolor{blue!70!black}{\textbf{—— 物理2·声光热 完 ——}}}

\vspace{0.3cm}
\textcolor{gray}{\small 下一本预告：物理3·电磁入门——从电路连接到达·芬奇没见过的电磁感应}
\end{center}

\newpage
\section*{\textcolor{defblue}{参考答案}}

\textbf{第一课 练一练：}
1. ×（真空不能传声）\quad
2. 振动，声带\quad
3. 340\quad
4. $s=vt=340\times4=1360\ \text{m}$\quad
5. A\quad
6. 固，气，340\quad
7. ×（还需要介质传播）\quad
8. 0.1，17

\textbf{第二课 练一练：}
1. 音调\quad
2. B\quad
3. 频率\quad
4. 范围\quad
5. $s=vt/2=340\times0.02/2=3.4\ \text{m}$\quad
6. B\quad
7. 赫兹（Hz），20 Hz，20000 Hz\quad
8. ×（音调与传播速度无关）

\textbf{第三课 练一练：}
1. B（月亮反射太阳光，不是光源）\quad
2. 同种均匀\quad
3. $3\times10^8$\quad
4. ×（光年是距离单位）\quad
5. $s=ct=3\times10^8\times500=1.5\times10^{11}\ \text{m}$\quad
6. C\quad
7. $3\times10^8$，慢\quad
8. $\surd$

\textbf{第四课 练一练：}
1. 入射角$=90^\circ-30^\circ=60^\circ$，反射角$=60^\circ$\quad
2. B\quad
3. $2+2=4\ \text{m}$\quad
4. C（换一面更大的镜子）\quad
5. 入射角$=90^\circ-40^\circ=50^\circ$，反射角$=50^\circ$\quad
6. B\quad
7. 遵循\quad
8. ×（平面镜成虚像，不能用光屏承接）

\textbf{第五课 练一练：}
1. 小于\quad
2. C\quad
3. 远离\quad
4. 下\quad
5. 是（空气是快介质，液体是慢介质）\quad
6. A（光从水到空气远离法线，人眼反向延长线看到更高的像）\quad
7. 小于，大于\quad
8. $\surd$

\textbf{第六课 练一练：}
1. 会聚，发散\quad
2. C（$u>2f$ → $20>2f$ → $f<10\ \text{cm}$）\quad
3. 凸，二倍焦距\quad
4. 凹，发散\quad
5. $u=40>2f$ → 倒立缩小的实像\quad
6. B\quad
7. 虚实，大小\quad
8. $\surd$

\textbf{第七课 练一练：}
1. ①熔化 ②液化 ③凝华\quad
2. 不变\quad
3. B\quad
4. 低\quad
5. $Q=cm\Delta t=2.1\times10^3\times0.5\times10=1.05\times10^4\ \text{J}$\quad
6. D\quad
7. 熔点\quad
8. $\surd$

\textbf{第八课 练一练：}
1. 做功，热传递\quad
2. $4.2\times10^3$\quad
3. $Q=cm\Delta t=0.88\times10^3\times0.5\times(75-25)=2.2\times10^4\ \text{J}$\quad
4. ×（热量是传递的能量，不是"含有"的）\quad
5. 高温，低温\quad
6. $c=Q/(m\Delta t)=1.68\times10^5/(2\times40)=2.1\times10^3\ \text{J/(kg·℃)}$\quad
7. C\quad
8. $4.6\times10^7$\quad
9. ×（热量是过程量，不能说"含有"）

\end{document}
